\section{Limitations and threats to validity}
\label{sec:limitations}

\begin{itemize}
  \item \textbf{The posterior is exact with respect to the rules, not to the
        opponents.} Constraints (C1)--(C5) are logical consequences of the
        rules and hold against any opponent whatsoever. The uniform prior over
        deals consistent with them is not, however, the full Bayesian posterior:
        that would additionally weight each deal by how likely the observed
        actions are under the opponents' actual policies, which is playstyle
        dependent. We implement and fit the natural correction \eqref{eq:prior}
        and find it adds nothing once the certificates are enforced, but that is
        evidence about this opponent population, not a theorem.
  \item \textbf{Population dependence.} Being best in this population is not
        being unexploitable. Our local best-response probe searches one policy
        class with one optimiser; a differently-shaped adversary could do better.
  \item \textbf{No equilibrium claim.} Fish with three-player teams and no
        communication channel calls for a team-maxmin equilibrium with
        correlation, not a Nash equilibrium, and we compute neither. Nothing here
        bounds exploitability.
  \item \textbf{Fitted on simulated opponents.} There are no expert human game
        logs and no independently authored Fish engines to test against. Every
        opponent in this study was written by us or ported from v0.3, which is a
        real limit on external validity.
  \item \textbf{The value function is linear.} Ridge-fitted on 405{,}000
        self-play decision points it reaches $R^2 = 0.29$ against the final
        half-suit differential
        (\path{research/v04/results/E14-valuefit-stats.txt}), which caps how much
        the one-ply lookahead and the stopping rule built on it can extract, and
        is consistent with the ablations being unable to separate either from
        zero.
  \item \textbf{Termination is imposed, not derived.} Theorem~\ref{thm:locked}
        makes patience correct, and \S\ref{sec:termination} shows that correct
        patience on both sides deadlocks. Our escalation rule terminates every
        game in this study, but its horizons are hand-set constants, not the
        output of any analysis, and a differently-tuned pair of policies could
        still find a cycle inside them.
  \item \textbf{The exploitability probe is a lower bound.} It searches one
        parametric policy class with one optimiser for a fixed budget. That it
        fails to beat FishBot v0.4 is evidence, not a bound; a differently-shaped
        adversary, or a larger budget, could succeed.
  \item \textbf{Ablation interaction.} Features are correlated; a term can look
        unnecessary because a correlated term substitutes for it. The paired
        design controls the deals, not the collinearity.
  \item \textbf{Rule dialect.} Results are for the 54-card, nine-half-suit form
        with wrong declarations awarded to the opponents and no prearranged
        conventions. \S\ref{sec:robustness} measures sensitivity to some of
        these, not all.
\end{itemize}
